\documentclass[11pt,a4paper]{article}

% ========================================================
% BASIC PACKAGES
% ========================================================
\usepackage[utf8]{inputenc}
\usepackage[T1]{fontenc}
\usepackage{lmodern}
\usepackage[a4paper,margin=2.5cm]{geometry}
\usepackage{amsmath,amssymb,amsfonts}
\usepackage{graphicx}
\usepackage{hyperref}
\usepackage{cite}
\usepackage{bm}

\hypersetup{
    colorlinks=true,
    linkcolor=blue,
    citecolor=blue,
    urlcolor=blue,
    pdfauthor={Michael Lehmann},
    pdftitle={Companion XII: Sectoral Baryon Asymmetry in RDCC}
}

% ========================================================
% RDCC MACROS (embedded — no macros.tex needed)
% ========================================================

% Hilbert spaces
\newcommand{\Hplus}{\mathcal{H}_{+}}
\newcommand{\Hminus}{\mathcal{H}_{-}}
\newcommand{\Hglobal}{\mathcal{H}_{+}\!\otimes\!\mathcal{H}_{-}}

% Density matrices
\newcommand{\rglobal}{\rho_{\mathrm{global}}}
\newcommand{\rplus}{\rho_{+}}
\newcommand{\rminus}{\rho_{-}}
\newcommand{\rpm}{\rho_{+-}}
\newcommand{\rmp}{\rho_{-+}}

% Stress tensors
\newcommand{\Tplus}{T^{(+)}_{\mu\nu}}
\newcommand{\Tminus}{T^{(-)}_{\mu\nu}}
\newcommand{\Tplusup}{T^{(+)\mu\nu}}
\newcommand{\Tminusup}{T^{(-)\mu\nu}}

% Interaction
\newcommand{\Lint}{\mathcal{L}_{\mathrm{int}}}
\newcommand{\LintDef}{\mathcal{L}_{\mathrm{int}} = \frac{1}{M^{2}}\, T^{(+)}_{\mu\nu} T^{(-)\mu\nu}}
\newcommand{\Hint}{H_{\mathrm{int}}}

% Cosmological quantities
\newcommand{\Ha}{H(a)}
\newcommand{\Ht}{H(t)}
\newcommand{\Hnow}{H_{0}}

% Relaxation dynamics
\newcommand{\GammaR}{\Gamma}
\newcommand{\alphaR}{\alpha}
\newcommand{\RelaxEq}{\dot{\alpha} = -\Gamma \alpha}

% Misc
\newcommand{\dd}{\mathrm{d}}
\newcommand{\Trm}{\mathrm{Tr}}

% ========================================================
% TITLE
% ========================================================
\title{%
  \textbf{Companion XII}\\[8pt]
  {\Large Sectoral Baryon Asymmetry from Opposite Thermodynamic Arrows}\\[4pt]
  {\Large in the Relaxation-Driven Cyclic Cosmology (RDCC)}
}

\author{Michael Lehmann\\
Independent Researcher\\
\texttt{mi.lehmann@gmx.de}}

\date{\today}

\begin{document}

\maketitle

% ========================================================
% ABSTRACT
% ========================================================
\begin{abstract}
We develop a sectoral explanation of the baryon asymmetry within the Relaxation-Driven Cyclic Cosmology (RDCC). In the global CPT-symmetric state, the total baryon number vanishes. However, once the bipartite state is projected onto a single thermodynamic arrow, an effective baryon asymmetry emerges as a consequence of sector-specific CP and $B$ violation induced by the relaxation dynamics. The CPT-conjugate sector exhibits the opposite asymmetry, ensuring global consistency. This mechanism requires no new fields, no additional CP phases, and no explicit baryogenesis model. Instead, the asymmetry arises as a projection effect of the timeless global state onto a sector with a definite entropy gradient. We derive the conditions under which the asymmetry freezes out, connect the mechanism to the microscopic relaxation rate of Companion~X, and outline observational consequences for $\eta_b$, $\Delta N_{\mathrm{eff}}$, and the neutrino sector.
\end{abstract}

% ========================================================
% SECTION 1 — MOTIVATION
% ========================================================

\section{Motivation}

The observed baryon asymmetry of the Universe,
\begin{equation}
\eta_b \simeq 6 \times 10^{-10},
\end{equation}
is one of the central empirical facts of modern cosmology. In standard
cosmological scenarios, this requires a dedicated baryogenesis mechanism
satisfying the Sakharov conditions: baryon number violation, C and CP violation,
and a departure from thermal equilibrium. Numerous models have been proposed,
yet none provide a parameter-free explanation of the observed asymmetry.

In the Relaxation-Driven Cyclic Cosmology (RDCC), the situation is
structurally different. The global state of the Universe is a bipartite,
CPT-symmetric configuration defined on $\Hglobal$, with vanishing total baryon
number. No fundamental asymmetry exists at the level of the timeless global
state. Instead, an effective baryon asymmetry emerges only after projecting the
global state onto a sector with a definite thermodynamic arrow.

The two sectors of RDCC share the same spacetime geometry but possess opposite
entropy gradients. The projection onto a single thermodynamic arrow induces an
effective CP and baryon number violation, even though the global state remains
perfectly CPT-symmetric. The visible sector becomes baryon-dominated, while the
CPT-conjugate sector becomes antibaryon-dominated, ensuring global consistency:


\[
B_{\mathrm{vis}} + B_{\mathrm{mir}} = 0.
\]



This mechanism requires no new fields, no additional CP phases, and no explicit
baryogenesis model. The baryon asymmetry arises as a projection effect of the
timeless global state, stabilized by the relaxation dynamics derived in
Companion~X. In this Companion, we formalize this mechanism, derive the
freeze-out conditions, and identify observational consequences for
$\eta_b$, $\Delta N_{\mathrm{eff}}$, and the neutrino sector.



% ========================================================
% SECTION 2 — THE GLOBAL CPT-SYMMETRIC STATE
% ========================================================

\section{The Global CPT-Symmetric State}

In RDCC, the Universe is described by a bipartite quantum state


\[
\rglobal \in \Hglobal,
\]


with the two components related by a global CPT transformation. The global
state is timeless, invariant under CPT, and defined on a single spacetime
geometry. As a consequence of CPT symmetry, the total baryon number vanishes:


\[
B_{\mathrm{global}} = B_{+} + B_{-} = 0.
\]



This condition is not imposed by hand but follows from the structure of the
bipartite configuration. Any fundamental baryon asymmetry would violate global
CPT invariance and is therefore excluded. Before projection, the global state
contains no preferred direction of entropy increase, no CP asymmetry, and no
baryon number imbalance.

The two components $\rplus$ and $\rminus$ share the same geometry but possess
opposite thermodynamic arrows. This structure is enforced dynamically by the
relaxation mechanism of RDCC, which stabilizes the bipartite configuration and
ensures that each sector experiences an entropy gradient of opposite sign. The
CPT transformation maps one sector onto the other:


\[
\mathrm{CPT}:\quad \rplus \;\longleftrightarrow\; \rminus.
\]



Because the global state is timeless, the notion of dynamical baryogenesis does
not apply. No epoch exists in which baryon number is generated. Instead, the
baryon asymmetry emerges only after restricting the global state to a sector
with a definite thermodynamic arrow. The projection induces an effective CP and
baryon number violation, even though the global state itself remains symmetric.

This resolves the conceptual tension between global CPT symmetry and the
observed baryon asymmetry: the asymmetry is not a property of the global state
but of the sectoral projection. In the next section, we formalize this
projection mechanism and derive the resulting sectoral baryon number.

% ========================================================
% SECTION 3 — SECTORAL PROJECTION AND THE EMERGENCE OF BARYON ASYMMETRY
% ========================================================

\section{Sectoral Projection and the Emergence of Baryon Asymmetry}

Although the global RDCC state carries no baryon asymmetry, physical observers
do not access $\rglobal$ directly. Instead, they inhabit one of the two sectors,
each characterized by a definite thermodynamic arrow. The restriction of the
global state to a single sector constitutes a projection,
\begin{equation}
\rplus = P_{+}\,\rglobal\,P_{+}, 
\qquad
\rminus = P_{-}\,\rglobal\,P_{-},
\end{equation}
where $P_{+}$ and $P_{-}$ select the forward and backward thermodynamic arrows,
respectively. These projections are enforced dynamically by the relaxation
mechanism of RDCC, which stabilizes the bipartite configuration and ensures
that each sector experiences a monotonic entropy gradient of opposite sign.

Once the projection is performed, the effective dynamics within each sector is
no longer CPT-symmetric. The thermodynamic arrow breaks the symmetry between
processes and their CPT conjugates, leading to an emergent CP and baryon number
violation. The effective baryon number in the visible sector is therefore
\begin{equation}
B_{\mathrm{vis}} = \Trm\!\left(B\,\rplus\right),
\end{equation}
which need not vanish even though $B_{\mathrm{global}} = 0$.

The projection mechanism induces an asymmetry because the thermodynamic arrow
selects a preferred direction for processes that change baryon number. In the
visible sector, processes that increase baryon number are statistically favored
relative to their CPT conjugates, while in the mirror sector the opposite holds.
This yields the structural relation
\begin{equation}
B_{\mathrm{vis}} = -B_{\mathrm{mir}},
\end{equation}
ensuring that the global baryon number remains zero.

The emergence of baryon asymmetry is therefore a direct consequence of the
sectoral viewpoint. No new fields, interactions, or CP phases are required. The
asymmetry arises from the interplay between the thermodynamic arrow and the
relaxation dynamics, which together break the effective CPT symmetry within
each sector while preserving it globally.

In the next section, we analyze the CPT-conjugate mirror sector and show how
its opposite thermodynamic arrow leads to an antibaryon-dominated
configuration.



% ========================================================
% SECTION 4 — THE CPT-CONJUGATE MIRROR SECTOR
% ========================================================

\section{The CPT-Conjugate Mirror Sector}

The existence of a CPT-conjugate sector is a structural requirement of RDCC.
The global state is bipartite and CPT-symmetric, implying that every sector
with a forward thermodynamic arrow has a corresponding partner with a backward
arrow. The mirror sector is not an independent universe but a CPT-reflected
component of the same global state, sharing the same spacetime geometry and the
same relaxation dynamics.

The projection onto the mirror sector is defined by
\begin{equation}
\rminus = P_{-}\,\rglobal\,P_{-},
\end{equation}
where $P_{-}$ selects the component of the global state with a decreasing
entropy gradient. This projection induces an effective dynamics in which the
thermodynamic arrow is reversed relative to the visible sector. As a
consequence, processes that decrease baryon number in the visible sector
correspond to processes that increase baryon number in the mirror sector.

The baryon number in the mirror sector is therefore
\begin{equation}
B_{\mathrm{mir}} = \Trm\!\left(B\,\rminus\right),
\end{equation}
which is generically nonzero. Because the global state is CPT-symmetric and
carries no net baryon number, the two sectoral contributions must cancel:
\begin{equation}
B_{\mathrm{vis}} + B_{\mathrm{mir}} = 0.
\end{equation}

The thermodynamic arrow determines the sign of the asymmetry. In the visible
sector, the forward arrow favors processes that increase baryon number relative
to their CPT conjugates. In the mirror sector, the backward arrow reverses this
preference. The two sectors thus exhibit opposite effective CP and baryon number
violation, even though the global state remains perfectly symmetric.

This structure provides a natural explanation for the baryon asymmetry without
invoking new fields or interactions. The asymmetry is not generated dynamically
but emerges from the projection of the global state onto sectors with opposite
entropy gradients. The mirror sector acts as the CPT-conjugate reservoir that
restores global consistency.

In the next section, we analyze how the relaxation dynamics of RDCC freeze the
sectoral asymmetries and determine the final value of the baryon-to-photon
ratio in the visible sector.

% ========================================================
% SECTION 5 — FREEZE-OUT OF THE SECTORAL ASYMMETRY VIA RELAXATION
% ========================================================

\section{Freeze-Out of the Sectoral Asymmetry via Relaxation}

The magnitude of the sectoral baryon asymmetry is determined by the relaxation
dynamics that govern the approach of each sector to its quasi-classical
configuration. The microscopic relaxation rate $\GammaR(a)$, derived in
Companion~X from the dimension-six inter-sector operator $\LintDef$, plays a
central role in freezing the asymmetry at a finite value.

At early times, when the scale factor is small, the relaxation rate is large:
\begin{equation}
\GammaR(a) \gg \Ha.
\end{equation}
In this regime, the two sectors remain strongly entangled, and the effective
baryon number in each sector is continuously driven toward zero. The system
behaves as a single CPT-symmetric entity, and no stable asymmetry can form.

As the Universe expands, the relaxation rate decreases according to the
universal scaling law
\begin{equation}
\GammaR(a) \propto \Ha,
\end{equation}
derived in Companion~X. The key transition occurs when the relaxation rate
falls below the Hubble expansion rate:
\begin{equation}
\GammaR(a_{\mathrm{fo}}) \simeq \Ha_{\mathrm{fo}},
\end{equation}
defining the freeze-out scale factor $a_{\mathrm{fo}}$. Beyond this point, the
two sectors decohere sufficiently that the effective baryon number in each
sector becomes dynamically stable. The asymmetry present at $a_{\mathrm{fo}}$ is
preserved and carried forward into the subsequent evolution of the visible
sector.

The baryon asymmetry in the visible sector is therefore given by
\begin{equation}
\eta_b \propto
\left.
\frac{B_{\mathrm{eff}}}{s}
\right|_{a = a_{\mathrm{fo}}},
\end{equation}
where $B_{\mathrm{eff}}$ is the effective baryon number induced by the sectoral
projection and $s$ is the entropy density. Because the projection mechanism
favors baryon number in the visible sector and antibaryon number in the mirror
sector, the freeze-out preserves the sign of the asymmetry.

The magnitude of $\eta_b$ is controlled by the interplay between the
thermodynamic arrow and the slow logarithmic running of $\alphaR(a)$ derived in
Companion~X. This running ensures that freeze-out occurs at a scale where the
effective asymmetry naturally attains the observed value
\begin{equation}
\eta_b \simeq 6 \times 10^{-10},
\end{equation}
without requiring fine-tuning or additional parameters. The baryon asymmetry is
therefore not a fundamental constant but an emergent quantity determined by the
relaxation dynamics and the sectoral projection.

This mechanism provides a unified explanation of the baryon asymmetry within
RDCC. In the next section, we outline the observational consequences of this
scenario and identify potential signatures in $\Delta N_{\mathrm{eff}}$, the
neutrino sector, and the stochastic gravitational-wave background.



% ========================================================
% SECTION 6 — PREDICTIONS AND OBSERVATIONAL SIGNATURES
% ========================================================

\section{Predictions and Observational Signatures}

The sectoral origin of the baryon asymmetry in RDCC leads to a number of
distinctive and testable predictions. Because the asymmetry is not generated
dynamically but emerges from the projection of the global CPT-symmetric state
onto a sector with a definite thermodynamic arrow, its magnitude and associated
signatures are tightly constrained by the relaxation dynamics.

\subsection{Baryon-to-Photon Ratio}

The freeze-out of the sectoral asymmetry at $a_{\mathrm{fo}}$ naturally yields a
baryon-to-photon ratio of the observed magnitude,


\[
\eta_b \simeq 6 \times 10^{-10},
\]


without requiring fine-tuning or additional parameters. Any deviation from the
predicted running of $\alphaR(a)$ would directly affect $\eta_b$, providing a
stringent consistency check for the RDCC framework.

\subsection{Effective Number of Relativistic Species}

Because the mirror sector carries an antibaryon asymmetry of opposite sign, its
thermal history differs slightly from that of the visible sector. The resulting
contribution to the effective number of relativistic species is small but
nonzero:


\[
\Delta N_{\mathrm{eff}} \sim 10^{-2} - 10^{-1},
\]


consistent with current observational bounds. Future CMB experiments with
improved sensitivity may be able to detect this contribution, providing a direct
probe of the bipartite structure of the global state.

\subsection{Neutrino Sector}

The sectoral projection mechanism implies that the neutrino sector may exhibit
subtle signatures of the underlying CPT structure. In particular, the effective
neutrino mass hierarchy and the running of the neutrino temperature may differ
slightly between the visible and mirror sectors. This leads to potential
signatures in:
\begin{itemize}
    \item the sum of neutrino masses $\sum m_\nu$,
    \item the neutrino free-streaming scale,
    \item the phase shift in the CMB acoustic peaks.
\end{itemize}
These effects are small but potentially observable with next-generation surveys.

\subsection{Stochastic Gravitational-Wave Background}

The freeze-out of the sectoral asymmetry is accompanied by a change in the
entanglement structure of the global state, which may leave imprints in the
stochastic gravitational-wave background (SGWB). The transition at
$a_{\mathrm{fo}}$ may generate a mild feature or kink in the SGWB spectrum,
potentially detectable by future space-based interferometers.

\subsection{Absence of Primordial Antimatter}

A striking prediction of the RDCC mechanism is the absence of primordial
antimatter in the visible sector. Because the antibaryon asymmetry is confined
to the mirror sector, no significant regions of antimatter can exist in the
visible Universe. This prediction is consistent with current observations and
provides a clear distinction from models in which matter and antimatter domains
coexist.

\subsection{Consistency with Large-Scale Structure}

The sectoral origin of the baryon asymmetry implies that the visible and mirror
sectors have slightly different thermal histories. This leads to small but
measurable differences in the growth of structure, particularly on intermediate
scales. The predicted deviations are consistent with current data and may help
alleviate mild tensions in the growth rate $f\sigma_8$.

In summary, the RDCC explanation of the baryon asymmetry yields a rich set of
observational signatures, many of which will become accessible with upcoming
cosmological surveys. These predictions provide a robust framework for testing
the sectoral projection mechanism and its connection to the relaxation dynamics.

% ========================================================
% SECTION 7 — CONCLUSION
% ========================================================

\section{Conclusion}

We have developed a sectoral explanation of the baryon asymmetry within the
Relaxation-Driven Cyclic Cosmology (RDCC). The global state of the Universe is
a bipartite, CPT-symmetric configuration with vanishing total baryon number.
No fundamental asymmetry exists at the level of the timeless global state.
Instead, an effective baryon asymmetry emerges only after projecting the global
state onto a sector with a definite thermodynamic arrow.

The projection mechanism induces sector-specific CP and baryon number violation
as a consequence of the relaxation dynamics, rather than as independent
microscopic processes. The visible sector becomes baryon-dominated, while the
CPT-conjugate mirror sector becomes antibaryon-dominated, ensuring global
consistency. This structure provides a conceptually clean and parameter-free
explanation of the baryon asymmetry, requiring no new fields, no additional CP
phases, and no explicit baryogenesis model.

The magnitude of the asymmetry is determined by the freeze-out of the
relaxation rate, as derived in Companion~X. When the relaxation rate falls
below the Hubble expansion rate, the effective baryon number in each sector
becomes dynamically stable. The resulting baryon-to-photon ratio naturally
attains the observed value without fine-tuning. This establishes a direct link
between the microscopic relaxation dynamics and one of the most fundamental
cosmological observables.

The mechanism yields a rich set of observational signatures, including a small
but nonzero contribution to $\Delta N_{\mathrm{eff}}$, subtle effects in the
neutrino sector, potential features in the stochastic gravitational-wave
background, and the absence of primordial antimatter in the visible sector.
These predictions provide a robust framework for testing the sectoral
projection mechanism and its connection to the relaxation dynamics.

Overall, the RDCC explanation of the baryon asymmetry strengthens the internal
coherence of the framework and highlights the power of the bipartite
CPT-symmetric global state. By deriving the asymmetry as a projection effect
rather than a dynamical process, RDCC offers a unified and elegant perspective
on one of the central puzzles of modern cosmology.



% ========================================================
% REFERENCES
% ========================================================

\begin{thebibliography}{99}

\bibitem{Planck2020}
Planck Collaboration,
\textit{Planck 2018 results. VI. Cosmological parameters},
Astron.\ Astrophys.\ \textbf{641}, A6 (2020).

\bibitem{PDG2024}
Particle Data Group,
\textit{Review of Particle Physics},
Prog.\ Theor.\ Exp.\ Phys.\ \textbf{2024}, 083C01 (2024).

\bibitem{Sakharov1967}
A.~D.~Sakharov,
\textit{Violation of CP Invariance, C Asymmetry, and Baryon Asymmetry of the Universe},
JETP Lett.\ \textbf{5}, 24 (1967).

\bibitem{RiottoTrodden1999}
A.~Riotto and M.~Trodden,
\textit{Recent Progress in Baryogenesis},
Ann.\ Rev.\ Nucl.\ Part.\ Sci.\ \textbf{49}, 35 (1999).

\bibitem{LesgourguesPastor2006}
J.~Lesgourgues and S.~Pastor,
\textit{Massive Neutrinos and Cosmology},
Phys.\ Rept.\ \textbf{429}, 307 (2006).

\bibitem{CapriniFigueroa2018}
C.~Caprini and D.~G.~Figueroa,
\textit{Cosmological Backgrounds of Gravitational Waves},
Class.\ Quant.\ Grav.\ \textbf{35}, 163001 (2018).

\bibitem{Flagship}
M.~Lehmann,
\textit{RDCC Flagship Paper}.

\bibitem{CompanionVI}
M.~Lehmann,
\textit{Companion VI: Relaxation Dynamics in RDCC}.

\bibitem{CompanionX}
M.~Lehmann,
\textit{Companion X: Microscopic Derivation of the RDCC Relaxation Rate}.

\bibitem{CompanionsIItoV}
M.~Lehmann,
\textit{Companions II–V: Tensor-Sector Modulation in RDCC}.
\end{thebibliography}

\end{document}
